% Gauss, Werke X/1 (1917), native60–62, printed53–55.
\SourceLeaf{060}\NoSourceHead
\SourceGroup{VIII}
\HistoricTitle{[ZUM REZIPROZITÄTSGESETZ DER QUADRATISCHEN UND DER BIQUADRATISCHEN RESTE.]}
\HistoricCitation{[Auf einem Deckelblatt in Ec2, Kapsel 43.]}
\NumberedItem{1}
\begin{center}\textit{Das Fundamentaltheorem für die quadratischen Reste.}\end{center}
\par\indent Es sei
\[
\begin{aligned}
\left(\frac ba\right)&=1&&\text{wenn }b\mathrm{R}a\\
\left(\frac ba\right)&=-1&&\text{wenn }b\mathrm{N}a.
\end{aligned}
\]
\noindent Dann ist ($a,b$ ungerade positiv)
\[
2\left(\frac ba\right)\left(\frac ab\right)=1+\left(\frac{-1}a\right)+\left(\frac{-1}b\right)-\left(\frac{-1}a\right)\left(\frac{-1}b\right).
\]
\NumberedItem{2}
\par\indent Die verschiedenen Zahlen, die für den Modulus $a+ib$ allen Zahlen congruent sind, lassen sich am besten in \textit{zwei} Complexen vorstellen; nemlich
\[
\begin{array}{rl}
1)&\alpha+i\alpha',\ \text{wo }\alpha,\alpha'\text{ alle geraden Zahlen zwischen }-a\text{ und }+a,\\
2)&\beta+i\beta',\ \text{wo }\beta,\beta'\text{ alle ungeraden Zahlen zwischen }-b\text{ und }+b
\end{array}
\]
\noindent vorstellen.
\clearpage

\SourceLeaf{061}\ArithmeticHead{54}{L}{ARITHMETIK.\quad NACHLASS.}
\NumberedItem{3}
\par\indent Sind $a+bi$, $a'+b'i$ Primzahlen; $a,a'$ ungerade, $b,b'$ gerade und
\[
\begin{aligned}
(a'+b'i)^{\tfrac14(aa+bb-1)}&\equiv i^\mu\ (\mathrm{mod.}\ a+bi)\\
(a+bi)^{\tfrac14(a'a'+b'b'-1)}&\equiv i^{\mu'}\ (\mathrm{mod.}\ a'+b'i),
\end{aligned}
\]
\noindent so ist
\[
\mu-\mu'-\tfrac12\{(a'-1)b+(a-1)b'+bb'\}
\]
\noindent durch $4$ theilbar.

Elegant in reellen Zahlen so:
\[
\begin{gathered}p=aa+bb\\\pi=\alpha\alpha+\beta\beta\end{gathered}
\qquad a\equiv\alpha\equiv1\ (\mathrm{mod.}\ 4)
\]
\[
\begin{aligned}
\text{Character }\frac{a\alpha+b\beta}{p}&=\text{Character }\frac{a\alpha+b\beta}{\pi}+\tfrac14(a\alpha+b\beta-1)\\
\text{Character }\frac{a\alpha-b\beta}{p}+\text{Character }\frac{a\alpha-b\beta}{\pi}&=\tfrac14(a\alpha-b\beta-1).
\end{aligned}
\]
\par\indent Noch schöner lässt sich das F[undamental] T[heorem] so generalisiren.

Es seien $m,m'$ ungerade komplexe Zahlen,\newline
\hspace*{6em}$p,p'$ ihre Normen,
\[
\begin{aligned}
m&\equiv i^\lambda\ (\mathrm{mod.}\ 2+2i)\\
m'&\equiv i^{\lambda'}\ (\mathrm{mod.}\ 2+2i);
\end{aligned}
\]
\noindent Dann ist
\[
\mathrm{Ch}\frac m{m'}-\mathrm{Ch}\frac{m'}m=\tfrac14(p'-1)\lambda-\tfrac14(p-1)\lambda'+\tfrac18(p-1)(p'-1).
\]
\begin{center}\textbf{BEMERKUNG.}\end{center}
\begingroup\small
\par\indent Unter $\mathrm{Ch}\frac m{m'}$, $\mathrm{Ch}\frac{m'}m$ ist hier der Exponent $\rho$ bezw. $\rho'$ in den Kongruenzen
\[
m^{\frac{p-1}{4}}\equiv i^\rho\ (\mathrm{mod.}\ m'),\quad m'^{\frac{p'-1}{4}}\equiv i^\rho\ (\mathrm{mod.}\ m)
\]
\noindent zu verstehen.
\begin{flushright}\textsc{Bachmann.}\end{flushright}
\clearpage\endgroup

\SourceLeaf{062}\ArithmeticHead{55}{R}{ZUM REZIPROZITÄTSGESETZ.}
\NumberedItem{4}
\HistoricCitation{[Aus Handbuch 21, Bg, S. 24.]}
\par\indent Das Fundamentaltheorem für die biquadratischen Reste ist
\[
\mathrm{Dec.}\frac{A+Bi}{a+bi}-\mathrm{Dec.}\frac{a+bi}{A+Bi}=\tfrac12(A-1)b+\tfrac12(a-1)B+\tfrac12Bb.
\]
\begin{center}\textbf{BEMERKUNG.}\end{center}
\begingroup\small
\par\indent Eine Datierung der Aufzeichnungen [1]—[3.] ist nicht gelungen. Für [4.] ergibt sich als untere Zeitgrenze das Ende des Jahres 1813.
\begin{flushright}\textsc{Bachmann.}\end{flushright}
\ClosingRule\clearpage\endgroup
